%======================================================================
\chapter{Methods}
%======================================================================

\section{Methods from Artificial Intelligence}

The technical architecture of Tobio rests on three interconnected computer vision tasks, each addressing a specific aspect of volleyball analysis. Rather than attempting to solve everything at once, the system was designed as a pipeline where each stage builds on the previous one. This modularity allowed for independent development and optimization of each component while maintaining flexibility for future expansion.

\subsection{Court Boundary Detection and Segmentation}

The first challenge was establishing a spatial reference frame. Before tracking anything else, the system needs to understand where the court is and how it maps to the video frame. This is more complex than it might seem—camera angles vary, lighting conditions change, and players frequently occlude court lines.

The solution employed YOLOv11, configured for instance segmentation rather than simple object detection. Where object detection provides bounding boxes, segmentation provides pixel-level masks that precisely outline the court boundaries. This distinction matters because volleyball courts have well-defined geometric properties that can be exploited once the boundaries are known.

The model was trained on a dataset of annotated volleyball court images sourced from Roboflow, which provided labeled segmentation masks for court boundaries under various viewing angles and lighting conditions. Training used mixed-precision optimization to accelerate convergence on available GPU hardware while maintaining numerical stability.

An important design decision was made during implementation: rather than tracking court position frame-by-frame, which would introduce unnecessary noise and computational overhead, the system processes the entire video to detect the court in frames where it is clearly visible, then averages the detected corner coordinates across all successful detections. This produces a single, stable set of four corner points that represent the court's position throughout the video. The assumption—that the camera remains stationary—holds for nearly all amateur volleyball recordings and dramatically simplifies downstream processing.

\subsection{Ball Trajectory Tracking}

With the court established, the next task was tracking the ball itself. A volleyball in flight is a small, fast-moving object that frequently becomes occluded by players or moves outside the frame. These gaps in detection create a fundamental problem: a tracker that only reports when it sees the ball will produce stuttering, incomplete trajectories.

The approach taken was two-staged. First, YOLOv11 was deployed again, this time configured for object detection to identify the ball in each frame. The model was trained on a separate dataset of annotated volleyball images where the ball was labeled with bounding boxes. During inference, the system processes each frame sequentially, recording the ball's position and confidence score whenever a detection exceeds the confidence threshold.

The second stage addresses the gaps. After initial detection, an interpolation algorithm examines the tracking data to identify missing frames. When a gap is found between two detections, the algorithm evaluates whether interpolation is appropriate by checking two conditions: the temporal gap must be small (fewer than 10 frames), and the spatial distance between the two detections must be reasonable (fewer than 150 pixels). If both conditions are met, linear interpolation fills in the missing frames by computing intermediate positions between the known detections.

This conditional interpolation strikes a balance between completeness and accuracy. It fills gaps caused by momentary occlusions or tracking failures while avoiding the creation of false trajectories when the ball genuinely moves out of frame or undergoes rapid, unpredictable motion.

\subsection{System Architecture and Integration}

The tracking components were integrated into a web-based platform built on a client-server architecture. The backend, implemented in Python using FastAPI, handles video processing and model inference. Models are stored remotely on AWS S3 and downloaded on-demand, allowing the system to be deployed without bundling large model files directly into the application.

The frontend, built with React and TypeScript, provides the user interface for uploading videos and visualizing results. When a video is uploaded, the backend processes it through the court and ball tracking pipelines in sequence, returning JSON-formatted tracking data. The frontend then renders this data as overlays on the video using HTML5 Canvas, allowing users to see detected court boundaries and ball trajectories synchronized with playback.

The decision to separate tracking from visualization—processing video server-side but rendering results client-side—was deliberate. It keeps the computationally intensive inference isolated to the backend where GPU resources are available, while allowing the lightweight visualization to run smoothly in any modern web browser.

\section{Datasets}

The effectiveness of any supervised learning system depends critically on the quality and relevance of its training data. Tobio relied on publicly available datasets sourced from Roboflow, a platform that aggregates and hosts computer vision datasets in standardized formats.

\subsection{Court Segmentation Dataset}

Court tracking used a dataset containing annotated volleyball court images with pixel-level segmentation masks. The dataset included images captured from various camera angles and distances, representing the diversity of perspectives encountered in real gameplay footage. Each image was labeled with a segmentation mask that precisely outlined the court boundaries, providing the ground truth needed for training an instance segmentation model.

The use of segmentation masks rather than simpler annotations (such as corner keypoints) was intentional. While keypoint annotation might seem more direct for a task that ultimately aims to extract court corners, segmentation provides richer spatial information that helps the model learn to distinguish court boundaries from other linear features in the frame, such as lines on walls or floors adjacent to the court.

\subsection{Ball Detection Dataset}

Ball tracking used a separate dataset of volleyball images annotated with bounding boxes around the ball. This dataset posed unique challenges: volleyballs are small relative to frame size, often partially occluded, and can appear blurred due to motion. The dataset needed to include examples of these difficult cases to train a model that could reliably detect the ball under realistic conditions.

The dataset was sourced from a Roboflow repository specifically curated for volleyball ball detection, containing hundreds of labeled examples across various game scenarios—serves, rallies, and out-of-play moments. This diversity was essential for generalizing beyond the controlled conditions of a training set to the messy reality of actual game footage.

\subsection{Data Augmentation and Preprocessing}

Both datasets underwent augmentation during training to improve model robustness. Standard augmentation techniques—rotation, scaling, brightness adjustment, and horizontal flipping—were applied to artificially expand the training set and reduce overfitting. These augmentations simulate the natural variability encountered in real footage, such as different lighting conditions and camera positions, without requiring additional labeled data.

\section{Validation Methods}

Validation in a project like this operates on two levels: technical validation of individual components and holistic validation of the integrated system. Each requires different metrics and approaches.

\subsection{Component-Level Validation}

For court segmentation, the primary metric was Intersection over Union (IoU), which measures the overlap between predicted segmentation masks and ground truth annotations. An IoU threshold of 0.5 was used to determine successful detections, a standard benchmark in semantic segmentation tasks. During training, the model's performance on a held-out validation set was monitored to detect overfitting and guide hyperparameter tuning.

Ball detection was evaluated using mean Average Precision (mAP), the standard metric for object detection tasks. This metric accounts for both localization accuracy (whether the predicted bounding box overlaps sufficiently with the ground truth) and classification confidence (whether the model correctly identifies the object as a ball). A detection was considered successful if it achieved an IoU of at least 0.5 with the ground truth bounding box and exceeded the confidence threshold.

The interpolation algorithm's effectiveness was harder to quantify with traditional metrics because there is no ground truth for interpolated frames—by definition, these are frames where no detection occurred. Instead, validation focused on qualitative assessment: viewing interpolated trajectories and verifying that they appeared smooth and physically plausible.

\subsection{System-Level Validation}

Beyond component metrics, the ultimate test of the system was its performance on real game footage—specifically, footage from the Carleton University volleyball team, the original motivation for the project. This footage represented the exact use case the system was designed to serve: single-camera recordings of indoor volleyball games under typical gym lighting conditions.

Validation on this footage was primarily qualitative. Videos were processed through the full pipeline, and the resulting visualizations were reviewed to assess whether the tracking data was accurate and useful. Specific failure modes were documented: frames where the ball was not detected despite being visible, false positives where non-ball objects were tracked, and cases where court boundaries were incorrectly estimated.

This qualitative validation revealed insights that pure metrics could not capture. For example, the system performed well during active rallies but struggled during timeouts when players stood on the court lines, occluding them. Similarly, ball tracking was reliable during serves and passes but less consistent during fast spikes where motion blur became severe. These observations informed iterative improvements to the system and highlighted areas for future work.